% from ega2-1-fr.tex
\oldpage[II]{64}

canoniquement à $X\times_Y Y'$, en posant
$X=\operatorname{Proj}(\mathcal{S})$
(\hyperref[II.3.5.3-fr]{3.5.3})~; composant avec
$\operatorname{Proj}(u^\sharp)$ la première projection
$p:X\times_Y Y'\to X$, on obtient donc un morphisme
$v:G(u^\sharp)\to X$, que nous désignerons par $\operatorname{Proj}(u)$, et
qui est tel que le diagramme
\[
  \xymatrix{
    G(u^\sharp)\ar[r]^{v}\ar[d]
    & X\ar[d]\\
    Y'\ar[r]_{g}
    & Y
  }
\]
soit commutatif.

En outre, pour tout $\mathcal{S}$-Module gradué quasi-cohérent $\mathcal{M}$,
on a un $v$-morphisme canonique
\[
\label{II.3.5.6.1-fr}
  \nu:\widetilde{\mathcal{M}}\longrightarrow
  (g^*(\mathcal{M})\otimes_{g^*(\mathcal{S})}\mathcal{S}')^{\sim}
  |G(u^\sharp).
\tag{3.5.6.1}
\]

En effet, $\nu^\sharp$ s'obtient en composant les homomorphismes
\[
  v^*(\widetilde{\mathcal{M}})
  =W^*(p^*(\widetilde{\mathcal{M}}))
  \longrightarrow W^*((g^*(\mathcal{M}))^{\sim})
  \longrightarrow
  (g^*(\mathcal{M})\otimes_{g^*(\mathcal{S})}\mathcal{S}')^{\sim}
  |G(u^\sharp)
\]
où la première flèche provient de l'isomorphisme
(\hyperref[II.3.5.3-fr]{3.5.3}) et la seconde est l'homomorphisme
(\hyperref[II.3.5.2-fr]{3.5.2, (ii)})~; lorsque $\mathcal{S}$ est engendrée
par $\mathcal{S}_1$, il résulte de (\hyperref[II.3.5.2-fr]{3.5.2}) que
$\nu^\sharp$ est un isomorphisme.

Comme cas particulier de (\hyperref[II.3.5.6.1-fr]{3.5.6.1}), on a, pour
tout $n\in\mathbf{Z}$, un $v$-morphisme canonique
\[
\label{II.3.5.6.2-fr}
  \nu:\mathcal{O}_X(n)\longrightarrow
  \mathcal{O}_{X'}(n)|G(u^\sharp).
\tag{3.5.6.2}
\]
\end{env}

\subsection{Sous-préschémas fermés d'un préschéma
$\operatorname{Proj}(\mathcal{S})$.}
\label{subsection:II.3.6-fr}

\begin{env}[3.6.1]
\label{II.3.6.1-fr}
Soient $Y$ un préschéma,
$\varphi:\mathcal{S}\to\mathcal{S}'$ un homomorphisme de
$\mathcal{O}_Y$-Algèbres graduées quasi-cohérentes, de degré $0$. On dit que
$\varphi$ est (\textbf{TN})-\emph{surjectif} (resp.
(\textbf{TN})-\emph{injectif}, (\textbf{TN})-\emph{bijectif}) s'il existe $n$
tel que, pour $k\geq n$, $\varphi_k:\mathcal{S}_k\to\mathcal{S}'_k$ soit
surjectif (resp. injectif, bijectif). Lorsqu'il en est ainsi on ramène l'étude
du morphisme
$\Phi:\operatorname{Proj}(\mathcal{S}')\to\operatorname{Proj}(\mathcal{S})$
correspondant au cas où $\varphi$ est surjectif (resp. injectif, bijectif).
Cela se démontre comme dans (\hyperref[II.2.9.1-fr]{2.9.1}) (qui en est le cas
particulier correspondant à $Y$ affine) en utilisant
(\hyperref[II.3.1.8-fr]{3.1.8}). Au lieu de dire que $\varphi$ est
(\textbf{TN})-bijectif, on dit aussi que c'est alors un
(\textbf{TN})-isomorphisme.
\end{env}

\begin{proposition}[3.6.2]
\label{II.3.6.2-fr}
Soient $Y$ un préschéma, $\mathcal{S}$ une $\mathcal{O}_Y$-Algèbre graduée
quasi-cohérente, et soit $X=\operatorname{Proj}(\mathcal{S})$.
\begin{enumerate}
  \item[\rm(i)] Si $\varphi:\mathcal{S}\to\mathcal{S}'$ est un
    homomorphisme (\textbf{TN})-surjectif de $\mathcal{O}_Y$-Algèbres
    graduées, le morphisme correspondant
    $\Phi=\operatorname{Proj}(\varphi)$
    (\hyperref[II.3.5.1-fr]{3.5.1}) est défini dans
    $\operatorname{Proj}(\mathcal{S}')$ tout entier et est une immersion
    fermée de $\operatorname{Proj}(\mathcal{S}')$ dans $X$. Si $\mathcal{J}$
    est le noyau de $\varphi$, le sous-préschéma fermé de $X$ associé à
    $\Phi$ est défini par le faisceau d'idéaux quasi-cohérent
    $\widetilde{\mathcal{J}}$ de $\mathcal{O}_X$.
  \item[\rm(ii)] Supposons en outre $\mathcal{S}_0=\mathcal{O}_Y$,
    $\mathcal{S}$ engendrée par $\mathcal{S}_1$ et $\mathcal{S}_1$ de type
    fini. Soit $X'$ un sous-préschéma fermé de
    $X=\operatorname{Proj}(\mathcal{S})$, défini par un faisceau
    quasi-cohérent $\mathcal{I}$ d'idéaux de $\mathcal{O}_X$. Soit
    $\mathcal{J}$

